% \documentclass[../main.tex]{subfiles}
% \begin{document}

% \begin{center}
%     \Large{\textbf{ABSTRACT}}\\
% \end{center}
% \vspace{1cm}
% Capacitive sensing technology has become a standard for touch-based interaction, yet most systems treat all finger contacts as uniform inputs. Recognizing and classifying specific finger types (e.g., thumb, index, or middle finger) from capacitive images can significantly enhance human-computer interaction and device ergonomics. However, implementing such recognition models on battery-constrained edge devices requires a strict balance between classification accuracy and power efficiency. This thesis investigates the application of Spiking Neural Networks (SNNs)—a bio-inspired, event-driven computing paradigm—as an energy-efficient solution for capacitive image-based finger recognition.

% The research begins by proposing specialized data exploration methods for the CapfingerId dataset to extract distinctive spatial features of different finger types. To establish a performance benchmark, a standard Convolutional Neural Network (CNN) is developed, followed by the design of an equivalent SNN architecture to allow for a direct comparison of spiking dynamics. Furthermore, the study explores the integration of transformer-based global attention through a Spiking Vision Transformer (Spiking ViT) model. The proposed architectures are rigorously evaluated by measuring their recognition accuracy across various finger categories and their corresponding energy footprints. \textcolor{red}{[Insert your results here: e.g., The results show that the Spiking ViT achieves an accuracy of X\% while consuming only Y mJ per inference]}. This study demonstrates that SNNs provide a viable, low-power alternative to traditional deep learning models for finger-type classification in next-generation interactive devices.
% \begin{flushright}
% \begin{tabular}{@{}c@{}}
% Student\\
% \textit{(Signature and full name)}
% \end{tabular}
% \end{flushright}

% \end{document}

\documentclass[../main.tex]{subfiles}
\begin{document}

\begin{center}
    \Large{\textbf{ABSTRACT}}\\
\end{center}
\vspace{1cm}

Capacitive sensing technology is widely used in touch-based interfaces, yet most systems treat all finger contacts as identical inputs. Recognizing specific finger types from capacitive images can improve user interaction and device usability, but implementing such functionality on resource-constrained devices requires models that are both accurate and energy-efficient.

This thesis investigates the use of Spiking Neural Networks (SNNs) for capacitive image-based finger recognition. Two SNN-based architectures are developed and evaluated. The first is a RANC-based SNN model, designed and trained entirely within the RANC framework. As a fully spiking architecture, it offers strong potential for deployment on neuromorphic and FPGA-based hardware while achieving recognition performance comparable to the CNN baseline of the CapFingerID dataset with reduced computational complexity.

The second architecture is a Spiking Vision Transformer (Spiking ViT), which combines spiking computation with transformer-based feature extraction. By integrating global attention mechanisms into a hybrid CNN-SNN framework, the proposed model achieves improved classification performance. Experimental results show that the Spiking ViT surpasses the original CapFingerID baseline by approximately 6\% in classification accuracy.

The results demonstrate that SNN-based models, particularly transformer-enhanced architectures, provide an effective and energy-efficient solution for finger-type recognition on next-generation embedded systems.
\begin{flushright}
\begin{tabular}{@{}c@{}}
\textbf{STUDENT}\\[2.5cm]
\textbf{Tran Minh Hien}
\end{tabular}
\end{flushright}

\end{document}